\documentclass{beamer}

\usetheme{Madrid}
\usecolortheme{crane}

\usepackage[utf8]{inputenc}
\usepackage{url}
\usepackage{textpos}
\usepackage{multirow}
\usepackage{graphicx}

\definecolor{darkgreen}{rgb}{0.0, 0.3, 0.13}

\newcommand{\bluehref}[2]{\href{#1}{\textcolor{blue}{#2}}}

\definecolor{emph}{RGB}{10,150,10}
\definecolor{contrast}{RGB}{0,0,120}
\definecolor{problem}{RGB}{120,0,0}
\definecolor{nice}{RGB}{10,150,10}

\usepackage{listings,bera}
\definecolor{sourcecode}{RGB}{0,0,100}
\definecolor{keywords}{RGB}{200,90,0}
\definecolor{comments}{RGB}{60,179,113}


\title{Automating computations: an introduction to workflows}
\author[Konrad HINSEN]{Konrad HINSEN}
\institute[CBM/SOLEIL]
          {Centre de Biophysique Moléculaire, Orléans, France\\
           and\\
           Synchrotron SOLEIL, Saint Aubin, France}
\date{}

\begin{document}

\begin{frame}
\titlepage
\end{frame}

\begin{frame}{Automation matters}

  \begin{itemize}
  \item Interactive work is great for exploration, but rarely reproducible.
  \item Even if you write down everything you do, you will make mistakes.
  \item And even if you don't make mistakes, you are never \textit{sure}
        there are no mistakes.
  \item \textbf{You know exactly what you have done only if you can re-run your computation.}
  \end{itemize}

\end{frame}

\begin{frame}{Programs, scripts, notebooks, workflows}

  \begin{itemize}
  \item Variants of the same idea: automation of computations
  \item Different terms and tools at different levels of organization
  \item Program: complex algorithms and/or data structures
  \item Script: short program, simple algorithms, no data abstraction
  \item Notebook: script with comments and output
  \item Workflow: coordination of multiple programs, multiple computers, multiple datasets
  \end{itemize}

\end{frame}

\begin{frame}{Overview}

  \begin{enumerate}
  \item Why workflows?
  \item Workflow managers: a diverse family of tools
  \item How to choose a workflow manager?
  \item The descendants of \texttt{Make}
  \item Snakemake: Make for scientific computing
  \item Case study: the incidence of influenza-like illness revisited
  \end{enumerate}

\end{frame}

\section*{Why workflows?}

\begin{frame}{}

  \begin{center}
  \Large Why workflows?
  \end{center}

\end{frame}

\begin{frame}{The limits of scripts and notebooks}

  \begin{itemize}
  \item Linear control flow: hard to read when too long
  \item Not well adapted to parameter scans and other loops
  \item Linear list of results: hard to inspect when too numerous
  \item All or nothing: no minimal update after change of an input
  \item Parallelization opportunities difficult to exploit
  \end{itemize}

\end{frame}

\begin{frame}{Tasks: the building blocks of workflows}

  \begin{itemize}
  \item Task = step in a computation
  \item Code + input data $\rightarrow$ output data
  \item Often implemented as a short script (bash, Python, R, ...)
  \item Input and output data stored in files
  \end{itemize}

\end{frame}

\begin{frame}{Workflows: execution plans for tasks}

  \begin{itemize}
  \item One task's output is another task's input
  \item Simple control structures, e.g. parameter scans
  \end{itemize}

  TODO: image of a workflow graph

\end{frame}

\begin{frame}{Workflow execution}

  \begin{itemize}
  \item Run tasks one by one or in parallel ...
  \item ... on one computer or distributed among several ...
  \item ... respecting the constraints of the workflow:
    \begin{itemize}
    \item Never run a task before its input data is ready
    \item Avoid running a task if its output is already available
    \end{itemize}
  \item Deal with failures (machine crash, broken network connection, ...)
  \end{itemize}

\end{frame}

\section*{Workflow managers: a diverse family of tools}

\begin{frame}{}

\end{frame}

\section*{How to choose a workflow manager?}

\begin{frame}{}

\end{frame}

\section*{The descendants of \texttt{Make}}

\begin{frame}{}

\end{frame}

\section*{Snakemake: Make for scientific computing}

\begin{frame}{}

\end{frame}

\section*{Case study: the incidence of influenza-like illness revisited}

\begin{frame}{}

\end{frame}


\end{document}
